% MAS1004 Assignment 2: long report.

\documentclass[11pt]{article}
\usepackage{mas1004}

\reporttitle{Assignment 2: Long report}
\studentname{김형석}
\studentid{20231020}

\begin{document}
\makeheader

\textbf{Final releases.} Track A: \texttt{track-a-1}, the starter recipe. Track B:
\texttt{track-b-2}, the \texttt{b2} run. Leaderboard scores were not available when
this report was prepared: my nickname did not appear on the leaderboard updated
2026-11-09 02:41 KST.

\section*{Track A}

\begin{center}\footnotesize
\begin{tabular}{lp{0.57\linewidth}r}
\toprule
tag & what changed from the run before & validation \\
\midrule
baseline (seed 0) & Starter recipe as it is; all training photos (44,257). Reference run. & 0.4984 \\
baseline (seed 1) & Only the seed changed (0 $\rightarrow$ 1), to see movement by chance. & 0.4979 \\
drop91 (seed 0) & Removed 91 photos from \texttt{selection.csv} (44,257 $\rightarrow$ 44,166). & 0.4972 \\
\bottomrule
\end{tabular}
\end{center}

The baseline was run to have a reference score for every later change. It gets
about half of the photos right (0.4984), well above the 20\% of random guessing;
Busan (0.375) and Osaka (0.346) were the weakest cities. Changing only the seed
moved the overall score by 0.0005, but single cities moved more. Therefore an
overall change under about 0.01 may be luck, and two seeds are too few to measure
the random variation reliably.

For \texttt{drop91}, the reason and interpretation are those in the lab notebook:
the score fell from 0.4984 to 0.4972 ($-0.0012$), far below the 0.01 limit of
luck, so I cannot say the change hurt or helped. The run cannot distinguish
between the hypothesis being wrong and too few photos being removed. The final
Track A release is therefore \texttt{track-a-1}, the starter recipe; \texttt{drop91}
did not beat baseline seed 0 and was not released. Its validation score is 0.4984.
Leaderboard score: not available.

\section*{Decisions you made from photos}

\textbf{No example mistakes figure is used.} The photos were inspected using
suspects sheets from the Track B \texttt{b1} model (0.4952).

I looked at training photos of all five cities by eye. I saw language written on
signs and building walls; traditional buildings, clothing and relics; food;
ports, fish markets and other scenes that cannot be seen in inland areas. About
20\% by eye were close-ups of leaves or animals such as sparrows and cats. About
30\% by eye were near-copies of the same scene from a slightly different angle.
Modern city scenes look mostly alike apart from regional and cultural clues such
as the sea, fish markets and language on signs.

The rule written before measuring was to remove training photos that carry no
linguistic or geographic clue, such as close-ups of leaves or animals, while
keeping photos with a clue even when the model got them wrong.

The change removed 91 photos. The count output is copied in
\texttt{outputs/curate\_count.txt}: 44,166 selected rows remain, with 11,254
Seoul, 5,179 Busan, 13,325 Tokyo, 7,183 Osaka and 7,225 Taipei. The validation
score changed from 0.4984 to 0.4972 ($-0.0012$). No extra photos were added.

I ran curate.py suspects and curate.py drop myself in PowerShell, following the rule I had written in LOG.md before measuring.

\section*{Track B}

The first valid Track B setup, \texttt{b1}, used ResNet18, 224px, 4 epochs and
\texttt{data/starter/train}; its validation score was 0.4952. The next run,
\texttt{b2}, added \texttt{--balance} and changed nothing else; its validation
score was 0.5107. The overall increase was 0.0155. Balancing draws every city
with equal probability (about 20\% each), instead of in proportion to photo
count, while the total amount of training stays the same.

I chose \texttt{b2} from the numbers, not from looking at photos. Busan rose by
0.075 and Tokyo fell by 0.127; Osaka rose by 0.123. Seoul fell by 0.060, One run cannot prove this.

I did not collect extra photos.

The final Track B validation score is 0.5107. Its leaderboard score was not
available, so a disagreement between validation and leaderboard scores cannot
be assessed.

The leaderboard uses frames from YouTube channels not in the starter data, and
the final set is phone street photos, so validation and leaderboard scores may disagree.

\subsection*{Your submission, as the grader sees it}
\outputfile{outputs/check_submission.txt}

\section*{What made the difference}

In Track A, \texttt{drop91} changed 0.4984 to 0.4972 ($-0.0012$), so the starter
recipe remained final. In Track B, \texttt{--balance} changed 0.4952 to 0.5107
($+0.0155$). \texttt{b2} was selected from the photo counts and per-city validation scores, not from looking at photos;
the log cautions that it was run only once.

With one more week, I would widen the suspects sheets (\texttt{--every 1}, larger
\texttt{--n}), remove more photos by the same rule, put Busan photos into
\texttt{selection.csv} twice, run \texttt{b2} with other seeds, inspect Osaka
mistakes, and repeat each change with seeds 0, 1 and 2.

\section*{All runs}
\outputfile{outputs/experiments.csv}

\section*{Lab notebook}
\outputfile{outputs/LOG.md}

\end{document}
